\scanpage{132}\mspage{67}

équivaut à la condition

(10)\qquad (To 3$'$)\qquad $Z^{\unsure{*|\sharp}}$ est un sous-groupe \deleted{\unsure{de $\mathfrak S_S$}}\inserted{$Z\subset\unsure{\mathfrak S|\mathcal S}_S$},

lequel sera donc nécessairement commutatif.

% slanted
\margin{\uline{NB} Il faudrait réexpliciter l'opération ``astucieuse'' de $Z$ sur $\Pi$ (p. 57, (35)), qui \uuline{n'est pas} le transport des \unsure{structure|structures}, et qu'il vaudrait donc mieux ne pas noter $i\mapsto\lambda i$, mais \unsure{avec notation|une notation} différente, genre $T_\lambda\cdot i$.}

On introduit maintenant
\[
\begin{aligned}
\Pi&\subset\unsure{\mathfrak P|\mathcal P|P}(\mathfrak P_2(S))\\[-2pt]
\rotatebox{90}{$=$}&\\[-2pt]
\{A_\lambda\mid\lambda\in Z\}&
\end{aligned}\tag*{(11)}
\]
\margin{Les groupes d'autom des $\Gamma_\lambda=(S,A_\lambda)$ sont tous égaux à $G_0$, et $Z$ opère sur $\Pi$ de façon transitive. $G_0$, $\Pi$ et $\Pi'$ ne changent pas si on remplace $(S,F_0)$ par $(S,F_i)$ $(i\in\Pi)$.}

\margin{\deleted{On avait \unsure{supposé}}\inserted{\unsure{étudié}} Pour la suite du \unsure{G|§}, on suppose $\Pi$ non réduit à un seul élément (sinon $G=G_0$, tout est trivial). Dans le cas de \unsure{(sous groupes)} la double pyramide, base $n$-gône, cela signifie que $n\equiv0\ (4)$ et $n\geqslant8$, \deleted{\unsure{auquel}}\inserted{en tous cas} dans le cas \unsure{bipyramidal|bipyramidale}, le \unsure{noyau|stabilisateur} de l'op. de $Z^*$ sur $\Pi$ est $\{1,\unsure{\underline{\circ}|\underline{-}|\varepsilon}\}$, \deleted{\illegible{1}}}

\[
\begin{aligned}
G&\subset\mathfrak S_S\\[-2pt]
\rotatebox{90}{$=$}&\\[-2pt]
\bigl\{f\in\mathfrak S_S\ \big|\ & f(A_0)\in\Pi\quad\deleted{\text{\unsure{ou encore|on}}}\ \inserted{\text{i.e. }\exists\lambda\in Z'\text{ t.q. }f(A_0)=A_\lambda}\\
&\text{i.e. } f(\Pi)=\Pi\bigr\}.
\end{aligned}\tag*{(12)}
\]

On \deleted{\unsure{fera|ferait}} note \deleted{que \unsure{les $F$|$A_i,F_i$} pour tout $\lambda\in Z$,}
\[
\deleted{\unsure{(13)}\quad\mathrm{To}\,4)}\qquad A_i\in\mathfrak P(\mathfrak P_2(S)),\quad F_i\in\mathfrak P(\mathfrak P_3(S)).\tag*{(13)}
\]
les ens. d'arêtes, resp. des faces, associés à un $i\in\Pi$ $(i=\lambda i_0,\deleted{\unsure{\forall}}\lambda\in Z)$. \inserted{il vaut mieux noter $T_\lambda i_0$} On \deleted{suppose que}\inserted{note que} $G$ ne dépend que de $\Pi\subset\mathfrak P(\mathfrak P_2(S))$ ou de l'ens. correspondant $\Pi'\subset\mathfrak P(\mathfrak P_3(S))$\unsure{;|.} On a

\diagram{\begin{tikzcd}[ampersand replacement=\&,row sep=small,column sep=small]
(14)\qquad\mathrm{Cent}(G)\subset Z\subset \& G \arrow[d,dash] \& \deleted{\text{et on}}\\
\deleted{\text{et on suppose}} \& G_0 \& {,}
\end{tikzcd}}
avec $G_0,Z$ distingués dans $G$, et on a
\[
G/G_0\lhook\joinrel\longrightarrow\mathrm{Aff}_Z(\Pi)\tag*{(15)}
\]
\margin{$[$\deleted{\unsure{Pour que}} Dans le cas bipyramidal. \unsure{p|P} avec $n\equiv0\ (4)$, $n\geqslant8$, il faut remplacer $Z$ par $Z/\{1,\underline{\circ}\}]$}
$G/G_0$ opère \uline{transitivement} sur $\Pi$

On suppose

% archive sheet number: 131
